\documentclass[11pt,a4paper]{article}

% Core font setup for a clean sans-serif document similar to the reference PDF.
\usepackage[T1]{fontenc}
\usepackage[utf8]{inputenc}
\usepackage{lmodern}

% Page size and margins. Adjust these if you want a tighter or looser layout.
\usepackage[a4paper,margin=22mm,headheight=32pt,headsep=18pt,footskip=18mm]{geometry}

% Table/layout helpers used by the template.
\usepackage{array}
\usepackage{fancyhdr}
\usepackage{graphicx}
\usepackage{longtable}
\usepackage{textcomp}

% Global paragraph and list styling.
\renewcommand{\familydefault}{\sfdefault}
\setlength{\parindent}{0pt}
\setlength{\parskip}{0.65em}
\setlength{\leftmargini}{1.5em}
\setlength{\leftmarginii}{1.2em}
\setlength{\itemsep}{0.35em}

% ---------------------------------------------------------------------------
% Document metadata
% Edit these values first when reusing the template.
% ---------------------------------------------------------------------------

% Brand or family label shown in the top-left header.
\newcommand{\brandname}{Neuromorpyx BrainBoard}

% Main product name shown on the cover and in the header.
\newcommand{\productname}{BrainBoard 15}

% Document type shown next to the product name.
\newcommand{\documenttype}{Datasheet}

% Secondary cover text under the title.
\newcommand{\documentsubtitle}{User Manual}

% Product code or SKU shown on the cover.
\newcommand{\sku}{SKU: BB15-0001}

% Footer date shown as the document modification date.
\newcommand{\documentmodified}{17/06/2026}

% Footer center text.
\newcommand{\documenttitlefull}{\productname\ \documenttype}

% Optional cover image path. Put the image next to this .tex file or change the path.
\newcommand{\coverimagepath}{./NMX-BB-15_F.jpg}

% Short market/application summary shown on the cover page.
\newcommand{\targetareas}{Embedded AI, prototyping, edge vision}

% ---------------------------------------------------------------------------
% Heading style
% This customizes section/subsection appearance without extra packages.
% ---------------------------------------------------------------------------

\makeatletter
\renewcommand\section{\@startsection{section}{1}{\z@}%
  {-3.5ex \@plus -1ex \@minus -.2ex}%
  {2.3ex \@plus.2ex}%
  {\normalfont\Large\bfseries\ttfamily}}
\renewcommand\subsection{\@startsection{subsection}{2}{\z@}%
  {-3.0ex \@plus -1ex \@minus -.2ex}%
  {1.4ex \@plus .2ex}%
  {\normalfont\large\bfseries}}
\renewcommand\subsubsection{\@startsection{subsubsection}{3}{\z@}%
  {-2.5ex \@plus -1ex \@minus -.2ex}%
  {1.0ex \@plus .2ex}%
  {\normalfont\normalsize\bfseries}}
\makeatother

% Rename the table of contents heading.
\renewcommand{\contentsname}{\ttfamily\bfseries CONTENTS}

% Simple boxed note macro.
% Usage:
%   \notebox{Title}{Body text goes here.}
\newcommand{\notebox}[2]{
  \begin{center}
  \setlength{\fboxsep}{8pt}
  \fbox{\parbox{0.94\linewidth}{\textbf{#1}\par\smallskip #2}}
  \end{center}
}

% ---------------------------------------------------------------------------
% Header/footer layout used on all pages.
% ---------------------------------------------------------------------------
\newcommand{\setbrandpagestyle}{
  \pagestyle{fancy}
  \fancyhf{}
  \fancyhead[L]{\ttfamily\bfseries\large \brandname}
  \fancyhead[R]{\ttfamily\bfseries\large \productname\\[-1pt]\documenttype}
  \renewcommand{\headrulewidth}{0.4pt}
  \fancyfoot[L]{\bfseries \thepage\ /\ \pageref{LastPage}}
  \fancyfoot[C]{\documenttitlefull}
  \fancyfoot[R]{Modified: \documentmodified}
}

% Cover image block.
% If the image file exists, it is shown.
% If it does not exist, a placeholder box is shown instead.
\newcommand{\coverimageblock}{
  \IfFileExists{\coverimagepath}{
    \includegraphics[width=0.31\textwidth]{\coverimagepath}
  }{
    \fbox{
      \begin{minipage}[c][45mm][c]{0.27\textwidth}
        \centering
        \ttfamily Optional cover image
      \end{minipage}
    }
  }
}

% ---------------------------------------------------------------------------
% Cover page
% Replace the placeholder description and target areas below.
% ---------------------------------------------------------------------------
\newcommand{\makedatasheetcover}{
  \thispagestyle{fancy}
  \vspace*{0.6cm}
  \begin{flushright}
    {\ttfamily\bfseries\LARGE \productname\\[2pt]\documenttype\par}
    \vspace{1.8cm}
    {\Large \documentsubtitle\par}
    \vspace{0.4cm}
    {\Large \sku\par}
  \end{flushright}

  \vspace{2.1cm}
  \begin{flushright}
    \coverimageblock
  \end{flushright}

  \vfill
  % Cover section: short product description.
  {\ttfamily\bfseries\Huge Description\par}
  \vspace{0.7cm}
  % Write 1 short paragraph here describing the board/product.
  Nicla style embeded module integrating the Brainchip AKD1500 Neuromorphic AI accelerator
  
  \vspace{1.0cm}
  % Cover section: target markets / use cases.
  {\ttfamily\bfseries\Huge Target Areas:\par}
  \vspace{0.55cm}
  % Write a short comma-separated list here.
  \targetareas

  \vspace{0.8cm}
  \rule{\textwidth}{0.5pt}
  \newpage
}

% Table of contents page.
\newcommand{\makedatasheettoc}{
  \thispagestyle{fancy}
  \tableofcontents
  \newpage
}

% Activate the shared header/footer style.
\setbrandpagestyle

\begin{document}

% Generate page 1 cover.
\makedatasheetcover

% Generate page 2 contents.
\makedatasheettoc

% ---------------------------------------------------------------------------
% Main body starts here.
% Replace each placeholder paragraph/table/list with your real content.
% ---------------------------------------------------------------------------

\section{Application Examples}
% Write a short introduction for the product use cases here.
Summarize the main use cases here. The reference document uses short introductory paragraphs
followed by compact bullet lists.

\begin{itemize}
  % Replace each bullet with a real application area.
  \item \textbf{Industrial automation:} Describe the primary industrial workflow or deployment.
  \item \textbf{Prototyping:} Explain how the board accelerates proof-of-concept work.
  \item \textbf{Edge AI:} Note the inference or sensing role in the wider system.
\end{itemize}

\section{Features}

\subsection{General Specifications Overview}
% Start with a short paragraph summarizing the product at a high level.
This section is usually a short paragraph plus a two-column summary table.

% Main spec summary table.
% Left column = spec label.
% Right column = value / explanation.
\begin{tabular}{|>{\bfseries}p{0.29\textwidth}|p{0.63\textwidth}|}
\hline
Feature & Description \\
\hline
Processor & Insert the MCU, accelerator, or SoC details here. \\
Memory & Flash, RAM, external storage, and relevant limits. \\
Connectivity & SPI, I2C, UART, USB, wireless, or other interfaces. \\
Power Supply & Supported voltage range and main power-input path. \\
Sensors & List onboard sensors and the purpose of each one. \\
Dimensions & Physical size in millimeters. \\
Weight & Board weight in grams. \\
\hline
\end{tabular}

\subsection{Accessories}
% Replace these with included or optional accessories.
\begin{itemize}
  \item Accessory or cable included with the product.
  \item Optional connector, header, or mechanical add-on.
\end{itemize}

\subsection{Related Products}
% Replace these with companion boards, modules, or kits.
\begin{itemize}
  \item Related base board or carrier.
  \item Related accessory or expansion module.
  \item Related development kit.
\end{itemize}

\section{Ratings}

\subsection{Recommended Operating Conditions}
% Use this table for electrical, thermal, or timing ranges.
Use a compact operating-range table for voltages, temperatures, and any required timing limits.

\begin{tabular}{|p{0.31\textwidth}|p{0.12\textwidth}|p{0.09\textwidth}|p{0.09\textwidth}|p{0.09\textwidth}|p{0.12\textwidth}|}
\hline
Parameter & Symbol & Min & Typ & Max & Unit \\
\hline
Input Voltage & $V_{IN}$ & 3.3 & 5.0 & 5.5 & V \\
Operating Temperature & $T_{OP}$ & -20 & 25 & 70 & \textdegree C \\
Current Draw & $I_{SYS}$ & -- & 120 & 250 & mA \\
\hline
\end{tabular}

% Use note boxes for practical warnings, power tips, or integration caveats.
\notebox{Power Tip}{Use short note boxes for practical guidance. The reference PDF uses these to
call out power and safety details without interrupting the flow of the section.}

\notebox{Safety Note}{Document any voltage-level caveats, thermal constraints, or handling
requirements here.}

\section{Functional Overview}

\subsection{Pinout}
% Put a pinout description here.
% If you have a pinout image, insert it with \includegraphics.
Place the pin summary here. If you have a diagram, insert it with \verb|\includegraphics|.

\subsection{Block Diagram}
% Put a system-level block diagram description or image here.
Add the architecture diagram or subsystem overview here.

\subsection{Power Supply}
% Explain rails, regulators, battery inputs, and power modes here.
Explain power domains, regulators, and any always-on versus switchable rails.

\section{Device Operation}

\subsection{Getting Started}
% Describe first-boot or first-use steps here.
Link the quick-start path here: setup, flashing, examples, and first validation steps.

\subsection{Libraries and Firmware}
% List required libraries, firmware, repos, and versions here.
List the core software package, repository URL, and any version constraints.

\subsection{Online Resources}
% Replace with actual docs, repos, tutorials, or product pages.
\begin{itemize}
  \item Documentation portal
  \item Library reference
  \item Example repository
\end{itemize}

\section{Mechanical Information}

\subsection{Board Dimensions}
% Put mechanical dimensions or an engineering drawing here.
Add the mechanical drawing or a dimensions table here.

\subsection{Connectors}
% Describe connector names, pitch, mating parts, and orientation here.
Describe connector types, mating parts, and mounting guidance.

\subsection{Peripherals and Actuators}
% Describe LEDs, buttons, microphones, sensors, and other onboard hardware here.
Summarize LEDs, buttons, microphones, cameras, or other user-facing hardware.

\section{Product Compliance}

\subsection{Compliance Summary}
% Keep this as a short summary table of regulatory compliance.
\begin{center}
\begin{tabular}{ll}
CE & European Union \\
RoHS & Restriction of Hazardous Substances \\
REACH & Registration, Evaluation, Authorization and Restriction of Chemicals \\
\end{tabular}
\end{center}

\subsection{Regional Statements}
% Put formal compliance declarations here.
Reserve this space for declarations, classifications, and labelling notes.

\section{Company Information}

% Company/contact table.
\begin{tabular}{|>{\bfseries}p{0.28\textwidth}|p{0.64\textwidth}|}
\hline
Company Information & Details \\
\hline
Company Name & BrainBoard Technologies \\
Company Address & Replace with your legal or support address. \\
\hline
\end{tabular}

\section{Reference Documentation}

% Reference list table.
% Replace the example URLs with real links, document names, and repositories.
\begin{longtable}{|p{0.08\textwidth}|p{0.34\textwidth}|p{0.48\textwidth}|}
\hline
No. & Reference & Link \\
\hline
\endhead
1 & Main documentation & \texttt{https://example.com/docs} \\
2 & Library repository & \texttt{https://example.com/library} \\
3 & Example projects & \texttt{https://example.com/examples} \\
\hline
\end{longtable}

\section{Document Revision History}

% Revision history table.
% Add one row per release or update.
\begin{tabular}{|p{0.18\textwidth}|p{0.14\textwidth}|p{0.58\textwidth}|}
\hline
Date & Revision & Changes \\
\hline
17/06/2026 & 1 & Initial template release. \\
\hline
\end{tabular}

% Label used by the footer page count.
\label{LastPage}

\end{document}
